\chapter{Campaña en 8 escenas: El laboratorio de B\&W}

\section{Estructura dramática}

La aventura está pensada para resolverse en 2 a 3 sesiones de juego de unas 3 a 4 horas cada una, con un ritmo cinematográfico \emph{in crescendo} que culmina en la medianoche del eclipse lunar del 8 de diciembre.

\section{Escena 1: El incendio en Babcock \& Wilcox}

\kobfigure{01_emergencia_bw.png}{Sirenas y humo negro saliendo de las instalaciones blindadas de Babcock \& Wilcox.}

Los personajes salen de clase en una tarde desapacible de otoño hacia el parque de Markonzaga. Mientras juegan o comentan las últimas novedades de los recreativos, una sacudida sorda hace vibrar los cristales de los edificios cercanos. A lo lejos, una columna de denso humo negro y destellos de un verde antinatural se elevan sobre las naves interiores de Babcock \& Wilcox.

Minutos después, camiones de bomberos y coches patrulla de la Policía Municipal y Protección Civil descienden por la Gran Vía hacia la vega de Galindo con las sirenas aullando.
\begin{itemize}
  \item \lead{Desafío (Astucia o Reflejos Dif. 5):} Si se acercan a los muros perimetrales, ven a los guardias de seguridad repeliendo a los periodistas locales y observan a una figura fría con bata blanca y gabardina: el \textbf{Dr. Victor Kessler}, increpando al jefe de seguridad: \emph{«¡El colector ha sufrido una purga y el sujeto se ha escapado! ¡Cierren las salidas del municipio ya!»}.
  \item \lead{Efectos ambientales:} Los walkman de los chavales emiten chillidos de estática, las agujas de los relojes digitales de cuarzo parpadean y un fuerte olor a ozono flota en el aire.
\end{itemize}

\section{Escena 2: El escondite secreto (La huerta)}

Tras cenar en sus casas (donde en el televisor Paco Lobatón en \emph{Quién sabe dónde} repasa el caso de niños desaparecidos en Euskadi y menciona el aniversario de la desaparición de Alex Etxebarria), el grupo se reúne a hurtadillas en su cuartel general: una caseta de aperos abandonada en una huerta en las afueras de Sestao.

Los jugadores comparten lo que saben sobre Alex, las actividades secretas de B\&W y las leyendas de las cuevas de La Arboleda. Planifican su primera excursión hacia los montes mineros para buscar pistas antes de que los adultos intervengan.

\section{Escena 3: La primera incursión en La Arboleda}

El grupo sube en bicicleta pedaleando con esfuerzo o cargando las bicis en el histórico funicular de La Reineta. En los montes de La Arboleda, el paisaje minero parece desierto bajo la niebla.

\begin{itemize}
  \item \lead{La bocamina acordonada:} Encuentran la boca de la cueva donde desapareció Alex cerrada con vallas de alambre de espino y carteles de \emph{«Propiedad privada: Babcock \& Wilcox»}.
  \item \lead{Desafío (Infiltración / Reflejos Dif. 7):} Sortear la valla sin cortar los cables de alarma de B\&W.
  \item \lead{Hallazgos:} En el interior de la galería no hay mineral de hierro corriente: las paredes muestran extrañas vetas cristalizadas que emiten luminiscencia verde. En el suelo descubren huellas de botas militares recientes y un cascabel de bicicleta oxidado que perteneció a Alex.
\end{itemize}

\section{Escena 4: La Sociedad de Ciencias Naturales}

\kobfigure{03_sociedad_ciencias.png}{Francisco en su despacho rodeado de fósiles de amonites y viejos informes.}

Buscando respuestas sobre los extraños cristales, el grupo acude a la Sociedad de Ciencias Naturales en el centro de Sestao para hablar con el paleontólogo \lead{Francisco}.

Francisco se muestra esquivo y asustado al principio, pero si los chavales demuestran que han visto las marcas de B\&W, se derrumba y les cuenta la verdad:
\begin{itemize}
  \item Les muestra los fósiles de amonites con aleaciones imposibles que desenterró en 1974.
  \item Confiesa que compartió sus notas con el Dr. Kessler y que este militarizó las excavaciones para extraer el Éter.
  \item Les advierte sobre el peligro mortal del \lead{eclipse del 8 de diciembre}: \emph{«Kessler cree que puede controlar la energía cósmica, pero si sobrecarga el colector, provocará una catástrofe que tragará toda la margen izquierda»}.
\end{itemize}

\section{Escena 5: La aparición de Alex}

\kobfigure{04_aparicion_alex.png}{Alex Etxebarria aparece en mitad de la noche empapado y rodeado de estática.}

Esa misma noche, mientras los chavales vigilan o duermen inquietos, unos golpes secos resuenan en la puerta de la caseta de la huerta. Al abrir, encuentran a \lead{Alex Etxebarria}.

Alex tiene once años, viste una bata raída de paciente de hospital bajo un chubasquero roto y sus ojos brillan con un halo esmeralda. Apenas recuerda su propio nombre ni el tiempo transcurrido, pero reconoce las voces de sus amigos. Objetos ligeros flotan a su alrededor cuando se emociona.
\begin{itemize}
  \item \lead{Mecánica de juego:} En este instante entra en juego la ficha de \textbf{Personaje Potenciado} y se reparten los aspectos de Alex entre los jugadores.
  \item \lead{La revelación de Alex:} Con voz entrecortada, Alex les advierte: \emph{«El hombre de la bata blanca me busca... quiere atarme a la máquina cuando la luna se vuelva roja»}.
\end{itemize}

\section{Escena 6: Emboscada en el parque de Markonzaga}

\kobfigure{05_emboscada_parque.png}{El Dr. Kessler y sus guardias jurados intentan capturar a Alex por la fuerza.}

Al día siguiente, mientras el grupo intenta cruzar el parque de Markonzaga para llevar a Alex a un refugio más seguro en casa de Francisco, son interceptados por dos furgonetas negras de las que descienden guardias de seguridad de B\&W armados con porras eléctricas y perros, liderados por el mismísimo \lead{Dr. Victor Kessler}.

\begin{itemize}
  \item \lead{El enfrentamiento:} Kessler exige la entrega inmediata del chico: \emph{«Ese niño no es humano, es un recurso estratégico del Estado. Apartaos si no queréis que vuestros padres pierdan el trabajo y terminen en la cárcel»}.
  \item \lead{Persecución en bicicletas (Reflejos / Combate Dif. 8):} Una trepidante persecución por las empinadas calles de Sestao, bajando escaleras y cruzando callejones estrechos donde las furgonetas no caben.
  \item \lead{Intervención de Alex:} Los jugadores pueden gastar Fichas de Adversidad para que Alex vuelque un contenedor con telequinesis o haga reventar las farolas para cegar a los perseguidores.
\end{itemize}

\section{Escena 7: El santuario de la colonia en La Arboleda}

\kobfigure{06_cueva_arboleda.png}{La nave alienígena biomecánica en las profundidades de la mina de La Arboleda.}

Comprendiendo que no podrán huir para siempre, Alex experimenta una intensa llamada psíquica y suplica a sus amigos volver a las profundidades de La Arboleda antes del eclipse.

En el corazón de la cueva, oculta tras una mampara derribada por B\&W, descubren la auténtica estructura de la \lead{nave colonia alienígena}: una inmensa cámara abovedada de roca y metal orgánico donde reposan cientos de monolitos fósiles con forma de conchas biomecánicas.
\begin{itemize}
  \item \lead{El glifo de resonancia:} Uno de los PJs descubre un glifo milenario en la roca. Al tocarlo con el colgante de Alex, la nave se ilumina con una frecuencia armónica que pacífica el Éter.
  \item \lead{La clave del clímax:} Francisco llega al lugar y comprende que para detener a Kessler no basta con apagar las máquinas de B\&W: deben redirigir la energía del colector central hacia el núcleo de la nave fósil para que se disipe sin peligro en la corteza terrestre.
\end{itemize}

\section{Escena 8: El clímax final en los laboratorios de B\&W}

\kobfigure{07_laboratorio_final.png}{Enfrentamiento final en el corazón del laboratorio de B\&W durante el eclipse lunar.}

\textbf{Noche del 8 de diciembre de 1992.} El cielo se despeja milagrosamente sobre la ría. La luna llena empieza a ser devorada por la sombra de la Tierra, tiñéndose de un color cobrizo y carmesí.

Guiados por Francisco y planos de ventilación conseguidos en la Sociedad de Ciencias, los personajes se infiltran en el búnker subterráneo de Babcock \& Wilcox en la vega de Galindo.
\begin{itemize}
  \item \lead{La cámara del Colector:} En una sala acorazada con compuertas de hormigón se alza el reactor de Éter: una gigantesca cúpula metálica conectada a turbinas y cables de alta tensión que emite un rugido electromagnético atronador.
  \item \lead{El choque final:} Kessler, respaldado por sus científicos y guardias restantes, intenta forzar el encendido del reactor. Los chavales deben usar su ingenio, tirar piedras con tirachinas a los paneles de control, cortar cables y coordinar a Alex para contener las descargas de plasma.
  \item \lead{La resolución:} Con el eclipse en su punto álgido, los jugadores ejecutan la inversión del colector guiados por Francisco y la mente de Alex. La energía es absorbida por la red terrestre de amonites, provocando un fogonazo de luz blanca que apaga los reactores y deja en silencio toda la factoría.
  \item \lead{Desenlace:} Kessler es detenido o desenmascarado ante la opinión pública; Francisco asume la custodia tutelar de Alex para ayudarle a reintegrarse en la vida escolar; y los chavales pedalean juntos al amanecer junto a la ría, sabiendo que han salvado a su pueblo de una catástrofe sin que ningún adulto se haya enterado jamás de lo ocurrido.
\end{itemize}
